\section*{Initial project description}

\textbf{Student}: Eske Haack, s214643 \\
\textbf{Graduate program}: M.Sc. in Human-Centered Artificial Intelligence \\
\textbf{Danish title}: Betingede Diffusionsmodeller til superopløsning af temperatur og nedbør over Europa \\
\textbf{English title}: Conditional Diffusion Models for Temperature and Precipitation Super-position over Europe \\
\textbf{ECTS Points}: 30 ECTS \\
\textbf{Starting Date}: 01-09-2026 \\
\textbf{ECTS Points during thesis}: 0 ECTS \\
\textbf{Company Name}: Danmarks Meteorologiske Institut \\
\textbf{Zip code (postnummer) for company}: 2100 København Ø \\
\textbf{Project carried out in}: Denmark \\
\textbf{Collaboration form}: Partly at DTU and partly at the company \\
\textbf{External supervisors}: Tian Tian, tian@dmi.dk, +45 24 81 32 03 \\
\textbf{Description: } This master thesis will investigate the application of conditional Diffusion Models on spatial super-resolution. Specifically the focus will be on downscaling daily temperature and precipitation data over the European Domain, using low-resolution EC-Earth3-Veg climate simulations as predictors and high-resolution HCLIM simulations as targets. The downscaled fields will be evaluated for their ability to represent future climate indicators, including heave-wave-relevant temperature extremes and heavy-precipitation characteristics.

The research aims to investigate if diffusion models are a viable option for temperature and precipitation downscaling, in a climate-projection setting. Diffusion models have been chosen for their generative properties, as well as the ability to compute ensembles, and thus uncertainties, of predictions.
In this study a series of diffusion-approaches and -architectures will be investigated to improve downscaling performance. As well as a shorter investigation on the models transferability and across scenarios and realizations. 